\documentclass[11pt]{article}
\usepackage[utf8]{inputenc}
\usepackage[T1]{fontenc}
\usepackage[spanish]{babel}
\usepackage{geometry,longtable,booktabs,array,enumitem,hyperref,xcolor,amssymb}
\geometry{margin=2.2cm}
\hypersetup{colorlinks=true,urlcolor=blue!60!black}
\title{Bitácora individual\\\large Hashing dinámico con hashing universal}
\author{\Nombre}
\date{Entrega: 11 de septiembre de 2026}
\begin{document}
\emergencystretch=3em
\maketitle

\section{Responsabilidad y criterio de trazabilidad}
Responsabilidad primaria propuesta: \textbf{\Rol}. Esta bitácora distingue evidencia verificable, borradores asistidos por IA y tareas planificadas. Una tarea pendiente sólo puede marcarse realizada cuando \Nombre{} la revise y adjunte evidencia coherente con el repositorio.

\section{Entradas}
\small
\begin{longtable}{p{1.7cm}p{2.1cm}p{3cm}p{3.2cm}p{4cm}}
\toprule Fecha/hito & Actividad & Decisión o aprendizaje & Evidencia & Estado/verificación \\
\midrule\endfirsthead
\toprule Fecha/hito & Actividad & Decisión & Evidencia & Estado \\\midrule\endhead
08-09-2026 / diagnóstico & Revisión del historial y asignación de responsabilidad & Preservar \texttt{main}; integrar selectivamente \texttt{gabo} & \EvidenciaInicial & \EstadoInicial \\
08-09-2026 / protocolo & Alineación con el diseño factorial de la investigación & Separar amortizado, esperado y promedio empírico & Commit \texttt{2efaf13} y documento de investigación & Integración generada con Codex; requiere revisión individual \\
08-09-2026 / piloto & Ejecución de 200 corridas y generación de ocho figuras & Conservar piloto separado de la corrida final & CSV piloto y ocho figuras & Artefactos generados; interpretación pendiente de revisión \\
09--10-09-2026 / revisión & \TareaAsignada & Registrar correcciones y poder explicarlas en defensa & \EvidenciaEsperada & \textbf{Pendiente}: completar por el estudiante \\
11-09-2026 / cierre & Reproducción y lectura final & No firmar resultados o fuentes no consultados & Commit o firma de validación & \textbf{Pendiente} \\
\bottomrule
\end{longtable}
\normalsize

\section{Uso de inteligencia artificial}
\subsection{Registro común de integración}
\begin{description}[leftmargin=3.4cm,style=nextline]
\item[Fecha y herramienta] 08-09-2026, ChatGPT/Codex.
\item[Propósito] Auditar ramas, alinear código con el protocolo, proponer pruebas, análisis, LaTeX y bitácoras trazables.
\item[Prompt representativo] ``Revisar \texttt{main} y \texttt{gabo}, integrar lo útil y crear documentos y correcciones conforme a las instrucciones y la investigación.''
\item[Fragmentos aprovechados] Instrumentación, pruebas, diseño factorial, estructura del artículo y anexos.
\item[Modificación/verificación] La suite automatizada pasó 30 pruebas. La revisión conceptual y bibliográfica debe completarla cada integrante.
\item[Responsabilidad] Codex no es fuente científica ni autor estudiantil; \Nombre{} responde sólo por los elementos que revise, comprenda y valide.
\end{description}

\subsection{Uso individual anterior}
\IAIndividual

\section{Lista de validación personal}
\begin{itemize}[label=$\square$]
\item Abrí y consulté las fuentes que aparecen asociadas a mi aporte.
\item Puedo explicar el algoritmo, las métricas y la diferencia entre amortizado y esperado.
\item Ejecuté o inspeccioné las pruebas relacionadas con mi tarea.
\item Corregí cualquier afirmación que no pudiera verificar.
\item Mi entrada coincide con commits, archivos y resultados del repositorio.
\end{itemize}

\vfill
\noindent Firma o confirmación del estudiante: \rule{7cm}{0.4pt}\\[1em]
Fecha: \rule{3cm}{0.4pt}
\end{document}
